\subsection{W-15 (TUR)}
Prediktivní a interpretativní metody vyhodnocování a testování uživatelských rozhraní. Testování použitelnosti.

\subsubsection*{UCD - User Centered Design}
\begin{itemize}
  \item Formativní fáze - analýza kontextu, testování použitelnosti
  \item Návrh - návrh a prototypování, architektura, logická struktura, vizuální návrh
  \item Sumativní fáze - závěrečné testování použitelnosti, pozorování artefaktu ve skutečném prostředí
\end{itemize}

\subsubsection*{Metody vyhodnocení}

\begin{itemize}
  \item Interpretativní metody - reální uživatelé, reálné prostředí, cíl je porozumění.
  \item Prediktivní metody - experti simulují uživatele, používá se prototyp, cíl je predikce a kontrola.
\end{itemize}

\subsubsection*{Testování použitelnosti}

\begin{itemize}
  \item Inspekce - heuristické, stránka po stránce se vyhodnotí, řeší se standardy, konzistence.
  \item Průchody - kognitivní, kontrola sekvencí, které vedou k určitým cílům
\end{itemize}

\subsubsection*{GOMS}
\begin{itemize}
  \item Goals - cíle uživatele
  \item Operators - manipulace s artefaktem - akce v aplikaci
  \item Methods - sekvence operátorů, které vedou k cíli
  \item Selection rules - výber operátorů nebo metod podle priorit/pravidel.
\end{itemize}
